\documentclass[12pt,a4paper]{article}

% ====== Packages ======
\usepackage[utf8]{inputenc}
\usepackage{amsmath, amssymb, amsfonts}
\usepackage{geometry}
\geometry{left=2.5cm, right=2.5cm, top=3cm, bottom=3cm}
\usepackage{xcolor}
\usepackage[most]{tcolorbox}
\usepackage{hyperref}

% ====== Colors Setup ======
\definecolor{mainblue}{RGB}{41, 128, 185}
\definecolor{maingreen}{RGB}{39, 174, 96}
\definecolor{mainred}{RGB}{192, 57, 43}
\definecolor{bgblue}{RGB}{235, 245, 251}
\definecolor{bggreen}{RGB}{233, 247, 239}
\definecolor{bgred}{RGB}{253, 237, 236}

% ====== TColorBox Environments ======
\newtcolorbox{defbox}[1][]{
    colback=bgblue,
    colframe=mainblue,
    fonttitle=\bfseries,
    boxrule=1pt,
    arc=4pt,
    enhanced,
    drop shadow,
    #1
}

\newtcolorbox{thmbox}[1][]{
    colback=bggreen,
    colframe=maingreen,
    fonttitle=\bfseries,
    boxrule=1pt,
    arc=4pt,
    enhanced,
    drop shadow,
    #1
}

\newtcolorbox{notebox}[1][]{
    colback=bgred,
    colframe=mainred,
    fonttitle=\bfseries,
    boxrule=1pt,
    arc=4pt,
    enhanced,
    drop shadow,
    #1
}

\title{\textbf{\color{mainblue} Probability and Statistics II Notes}}
\author{Zhiyao (Rachel) Lei}
\date{}

\begin{document}

\maketitle
\tableofcontents
\newpage

% Set section counter to start from 5 to match the notes chapters
\setcounter{section}{4}

% ==========================================
% Section 5: Multivariate Distribution
% ==========================================
\section{Multivariate Distribution}

\subsection{Joint Distributions}
\begin{itemize}
    \item \textbf{Discrete:} The joint probability mass function is
    \[
        p_{X,Y}(x,y)=P(X=x,Y=y).
    \]
    The marginal PMF and joint CDF are
    \[
        p_X(x)=\sum_y p_{X,Y}(x,y),
        \qquad
        F_{X,Y}(x,y)=P(X\le x,Y\le y).
    \]
    \item \textbf{Continuous:} A joint density satisfies
    \[
        \int_{-\infty}^{\infty}\int_{-\infty}^{\infty}
        f_{X,Y}(x,y)\,dy\,dx=1.
    \]
    Its marginal density is
    \[
        f_X(x)=\int_{-\infty}^{\infty}f_{X,Y}(x,y)\,dy.
    \]
\end{itemize}

\subsection{Expectation, Covariance, and Independence}
\begin{itemize}
    \item \textbf{Expected values:}
    \[
        E[g(X,Y)]
        =\int_{-\infty}^{\infty}\int_{-\infty}^{\infty}
        g(x,y)f_{X,Y}(x,y)\,dy\,dx.
    \]
    Linearity gives $E[aX+bY+c]=aE[X]+bE[Y]+c$.
    \item \textbf{Independence:} Random variables $X$ and $Y$ are independent if
    \[
        f_{X,Y}(x,y)=f_X(x)f_Y(y)
        \qquad\text{for all }x,y.
    \]
    Independence implies
    \[
        E[g(X)h(Y)]=E[g(X)]E[h(Y)]
    \]
    whenever the expectations exist. It also implies
    \[
        \operatorname{Var}(X+Y)
        =\operatorname{Var}(X)+\operatorname{Var}(Y).
    \]
    \item \textbf{Covariance:}
    \[
        \operatorname{Cov}(X,Y)
        =E[(X-\mu_X)(Y-\mu_Y)]
        =E[XY]-E[X]E[Y].
    \]
    \item \textbf{Correlation coefficient:}
    \[
        \rho_{X,Y}
        =\frac{\operatorname{Cov}(X,Y)}{\sigma_X\sigma_Y},
        \qquad \sigma_X,\sigma_Y>0.
    \]
\end{itemize}

\begin{defbox}[title=Conditional Distributions]
The conditional density of $Y$ given $X=x$ is
\[
    f_{Y\mid X}(y\mid x)=\frac{f_{X,Y}(x,y)}{f_X(x)}
\]
for values of $x$ such that $f_X(x)>0$. The conditional variance is
\[
    \operatorname{Var}(Y\mid X=x)
    =E(Y^2\mid X=x)-[E(Y\mid X=x)]^2.
\]
\end{defbox}


% ==========================================
% Section 6: Functions of Multivariate Random Variables
% ==========================================
\section{Functions of Multivariate Random Variables}

\subsection{Transformation Methods}
\begin{itemize}
    \item \textbf{CDF method:} For $W=h(X,Y)$,
    \[
        F_W(w)=P(h(X,Y)\le w)
        =\iint_{\{(x,y):\,h(x,y)\le w\}}
        f_{X,Y}(x,y)\,dx\,dy.
    \]
    \item \textbf{Jacobian method:} Let $(U,V)=T(X,Y)$ be one-to-one,
    with inverse $(X,Y)=T^{-1}(U,V)$. Then
    \[
        f_{U,V}(u,v)
        =f_{X,Y}\!\left(x(u,v),y(u,v)\right)
        \left|
            \frac{\partial(x,y)}{\partial(u,v)}
        \right|.
    \]
    \item \textbf{MGF method:} The joint moment-generating function is
    \[
        M_{U,V}(t_1,t_2)=E\!\left[e^{t_1U+t_2V}\right],
    \]
    when it exists in a neighborhood of $(0,0)$.
\end{itemize}

\subsection{Distributions of Sums \& Sample Statistics}
\begin{itemize}
    \item \textbf{Sums:}
    \[
        E\!\left[\sum_{i=1}^kX_i\right]
        =\sum_{i=1}^kE[X_i].
    \]
    If the variables are independent, then
    \[
        \operatorname{Var}\!\left(\sum_{i=1}^kX_i\right)
        =\sum_{i=1}^k\operatorname{Var}(X_i).
    \]
    \item \textbf{Convolution formula:} If $X$ and $Y$ are independent and
    $S=X+Y$, then
    \[
        f_S(s)=\int_{-\infty}^{\infty}f_X(x)f_Y(s-x)\,dx.
    \]
\end{itemize}

\begin{thmbox}[title=Sample Distributions and Order Statistics]
\textbf{Sample Statistics:} 
\begin{itemize}
    \item Sample mean:
    $\bar{X}=\frac{1}{n}\sum_{i=1}^nX_i$, with
    $E[\bar{X}]=\mu$ and $\operatorname{Var}(\bar{X})=\sigma^2/n$.
    \item Sample variance:
    $S^2=\frac{1}{n-1}\sum_{i=1}^n(X_i-\bar{X})^2$.
\end{itemize}
\textbf{Order statistics:} Write
$X_{(1)}=\min(X_1,\ldots,X_n)$ and
$X_{(n)}=\max(X_1,\ldots,X_n)$. For an i.i.d.\ continuous sample, the density
of the $k$-th order statistic is
\[
    f_{X_{(k)}}(x)
    =\frac{n!}{(k-1)!(n-k)!}
    [F(x)]^{k-1}[1-F(x)]^{n-k}f(x).
\]
\end{thmbox}


% ==========================================
% Section 7: Estimation
% ==========================================
\section{Estimation}

\subsection{Point Estimation}
\begin{itemize}
    \item A point estimator is a statistic used to estimate an unknown
    population parameter.
    \item The sample variance $S^2$ is an unbiased estimator of $\sigma^2$:
    $E[S^2]=\sigma^2$.
    \item If $X\sim\operatorname{Binomial}(n,p)$, then
    $\hat{p}=X/n$ satisfies
    \[
        E[\hat{p}]=p,
        \qquad
        \operatorname{Var}(\hat{p})=\frac{p(1-p)}{n}.
    \]
\end{itemize}

\subsection{Confidence Intervals (CI)}
\begin{defbox}[title={$100(1-\alpha)\%$ Confidence Intervals}]
\textbf{For a single population:}
\begin{itemize}
    \item \textbf{$\mu$ ($\sigma$ known):}
    $\bar{X}\pm z_{1-\alpha/2}\frac{\sigma}{\sqrt{n}}$.
    \item \textbf{$\mu$ ($\sigma$ unknown, normal population):}
    $\bar{X}\pm t_{1-\alpha/2,n-1}\frac{S}{\sqrt{n}}$.
    \item \textbf{Proportion $p$ (large sample):}
    $\hat{p}\pm z_{1-\alpha/2}
    \sqrt{\frac{\hat{p}(1-\hat{p})}{n}}$.
    \item \textbf{Variance $\sigma^2$ (normal population):}
    \[
        \left[
            \frac{(n-1)S^2}{\chi^2_{1-\alpha/2,n-1}},
            \frac{(n-1)S^2}{\chi^2_{\alpha/2,n-1}}
        \right].
    \]
\end{itemize}
\end{defbox}


% ==========================================
% Section 8: Hypothesis Testing
% ==========================================
\section{Hypothesis Testing}

\subsection{Fundamental Concepts}
\begin{itemize}
    \item \textbf{Steps:} Formulate $H_0$ and $H_1$, choose a test statistic,
    verify the assumptions, select a significance level, and determine the
    rejection region or p-value.
    \item \textbf{Errors and power:}
        \begin{itemize}
            \item Type I error:
            $P(\text{reject }H_0\mid H_0\text{ is true})=\alpha$.
            \item Type II error:
            $P(\text{fail to reject }H_0\mid H_0\text{ is false})=\beta$.
            \item Power:
            $P(\text{reject }H_0\mid H_0\text{ is false})=1-\beta$.
        \end{itemize}
    \item \textbf{p-value:} Under $H_0$, the p-value is the probability of a
    test statistic at least as extreme as the observed value. Reject $H_0$
    when the p-value is at most $\alpha$.
\end{itemize}

\begin{thmbox}[title=Tests for One and Two Populations]
\textbf{Single Population Tests:}
\begin{itemize}
    \item For $\mu$ ($\sigma$ known):
    $Z=\frac{\bar{X}-\mu_0}{\sigma/\sqrt{n}}\sim N(0,1)$ under $H_0$.
    \item For $\mu$ ($\sigma$ unknown, normal population):
    $T=\frac{\bar{X}-\mu_0}{S/\sqrt{n}}\sim t_{n-1}$ under $H_0$.
    \item For $p$ (large sample):
    $Z=\frac{\hat{p}-p_0}{\sqrt{p_0(1-p_0)/n}}
    \mathrel{\dot{\sim}}N(0,1)$ under $H_0$.
\end{itemize}
\textbf{Two Population Tests:}
\begin{itemize}
    \item Difference in means with known variances:
    \[
        Z=\frac{\bar{X}_1-\bar{X}_2-\Delta_0}
        {\sqrt{\sigma_1^2/n_1+\sigma_2^2/n_2}}.
    \]
    \item Ratio of variances for independent normal samples:
    \[
        F=\frac{S_1^2/\sigma_1^2}{S_2^2/\sigma_2^2}
        \sim F_{n_1-1,n_2-1}.
    \]
\end{itemize}
\end{thmbox}

\end{document}
